%% kapa: exam=1 round=21 year=2010 period=1 form=A subject=economics subject_ko=경제원론 seq=34 number=74
%% source: Q-Net 감정평가사 기출문제 (cst00309) artlSeq=5024181 file=extracted/1008/5024181/2010년 제21회 감정평가사 1차시험 기출문제/1교시.hwp
%% conversion: kordoc 4.24.51 -> tools/gen1.py (auto); review=proofread-sonnet; flags=none
\begin{question}{74}
독점기업 A는 완전가격차별을 하고 있다. 이에 관한 설명으로 옳은 것을 모두 고른 것은?
\begin{framedtext}
ㄱ. 총잉여(total surplus, net social benefit)는 가격차별을 하지 않을 때보다 적다. \par
ㄴ. 총잉여가 A기업에게 귀속된다. \par
ㄷ. 생산량은 가격차별을 하지 않을 때보다 많다. \par
ㄹ. A기업의 생산량은 동일한 시장수요와 비용을 지닌 완전경쟁시장에 비해 더 적다. \par
\end{framedtext}
\begin{choices}
  \item ㄱ, ㄴ
  \item ㄴ, ㄷ
  \item ㄷ, ㄹ
  \item ㄱ, ㄴ, ㄷ
  \item ㄴ, ㄷ, ㄹ
\end{choices}
\end{question}
\answer{2}
